% !TEX TS-program = lualatexmk
% !TEX parameter =  --shell-escape

\documentclass[vespers_book_la_eng.tex]{subfiles}

\ifcsname preamble@file\endcsname
  \setcounter{page}{\getpagerefnumber{M-vespers_book_la_eng_sundays_pp}}
\fi

\begin{document}

%\thispagestyle{empty}

\phantomsection
\chapter*{SUNDAYS AFTER PENTECOST.}\header{sundays after pentecost.}
\addcontentsline{toc}{chapter}{Sundays after Pentecost.}

\rubrique{All as in the psalter of Sunday, \normaltext{\pageref{M-vespers_book_la_eng_sunday_psalms},} except as given below.}

\bigtitle{IV Sunday after Pentecost.} 
\addcontentsline{toc}{section}{IV Sunday after Pentecost} \phantomsection

\gscore[Ad Magnif.]{Ant. 1. g}{an_praeceptor_solesmes_1961}
\antiphona{english}{Master, we have toiled all night, and have taken nothing; nevertheless at thy word I will let down the net}
%\textes{versiculus_pa}

\collectlateng

\texteparacol{\lettrine{D}{a} nobis, quǽsumus Dómine,~† ut et mundi cursus pacífice nobis tuo órdine dirigátur:~* et Ecclésia tua tranquílla devotióne lætétur. Per Dóminum.}
{english}{\noindent 
Grant, O Lord, we beseech thee, that the course of this world may be so peaceably ordered by thy governance, that thy Church may joyfully serve thee in all godly quietness.
Through our Lord.}

\bigtitle{V Sunday after Pentecost.}
\addcontentsline{toc}{section}{V Sunday after Pentecost} \phantomsection

\gscore[Ad Magnif.]{Ant. 1. D}{an_si_offers_munus_solesmes_1961}
\antiphona{english}{If thou bring thy gift to the Altar, and there rememberest that thy brother hath aught against thee, leave there thy gift before the Altar, and go thy way; first be reconciled to thy brother, and then come and offer thy gift. Alleluia}

\collectlateng

\texteparacol{\lettrine{D}{e}us, qui diligéntibus te bona invisibília præparásti:~† infúnde córdibus nostris tui amóris afféctum; ut te in ómnibus et super ómnia diligéntes, promissiónes tuas, quæ omne desidérium súperant, consequámur. Per Dóminum.}
{english}{\noindent 
O God, who hast prepared for them that love thee such good things as pass man's understanding; pour into our hearts such love toward thee, that we, loving thee in all things and above all things, may obtain thy promises, which exceed all that we can desire.
Through our Lord.}

\bigtitle{VI Sunday after Pentecost.}
\addcontentsline{toc}{section}{VI Sunday after Pentecost} \phantomsection

\gscore[Ad Magnif.]{Ant. 7. b}{an_misereor_super_turbam_solesmes_1961}
\antiphona{english}{I have compassion on the multitude, because they have now been with Me three days, and have nothing to eat; and if I send them away fasting, they will faint by the way. Alleluia}

\collectlateng

\texteparacol{\lettrine{D}{e}us virtútum, cujus est totum quod est óptimum:~† ínsere pectóribus nostris amórem tui nóminis, et præsta in nobis religiónis augméntum;~* ut quæ sunt bona nútrias, ac pietátis stúdio, quæ sunt nutríta, custódias.
Per Dóminum.}{english}
{\noindent O God of all power and might, who art the author and giver of all good things, graft in our hearts the love of thy Name, increase in us true religion, nourish us with all goodness, and, of thy great mercy, keep us in the same.
Through our Lord.}

\bigtitle{VII Sunday after Pentecost.}
\addcontentsline{toc}{section}{VII Sunday after Pentecost} \phantomsection

\gscore[Ad Magnif.]{Ant. 1. D}{an_non_potest_arbor_bona_solesmes_1961}
\antiphona{english}{A good tree cannot bring forth evil fruit, neither can a corrupt tree bring forth good fruit. Every tree that bringeth not forth good fruit, is hewn down and cast into the fire. Alleluia}

\collectlateng

\texteparacol{\lettrine{D}{eus}, cujus providéntia in sui dispositióne non fállitur:~† te súpplices exorámus; ut nóxia cuncta submóveas,~* et ómnia nobis profutúra concédas.
Per Dóminum.}{english}
{\noindent O God, whose never-failing providence ordereth all things both in heaven and earth, we humbly beseech thee to put away from us all hurtful things, and to give us those things which be profitable for us.
Through our Lord.}

\bigtitle{VIII Sunday after Pentecost.}
\addcontentsline{toc}{section}{VIII Sunday after Pentecost} \phantomsection

\gscore[Ad Magnif.]{Ant. 4. E}{an_quid_faciam_solesmes_1961}
\antiphona{english}{What shall I do? For my lord taketh away from me the stewardship. I cannot dig, to beg I am ashamed. I am resolved what to do, that, when I am put out of the stewardship, they may receive me into their houses}

\newpage

\collectlateng

\texteparacol{\lettrine{L}{a}rgíre nobis, quǽsumus Dómine, semper spíritum cogitándi quæ recta sunt, propítius et agéndi:~† ut qui sine te esse non póssumus,~* secúndum te vívere valeámus. Per Dóminum.}{english}
{\noindent Grant to us, Lord, we beseech thee, the spirit to think and do always such things as be rightful, that we, who cannot do anything that is good without thee, may by thee be enabled to live according to thy will.
Through our Lord.}

\bigtitle{IX Sunday after Pentecost.}
\addcontentsline{toc}{section}{IX Sunday after Pentecost} \phantomsection

\gscore[Ad Magnif.]{Ant. 8. G}{an_scriptum_est_enim_quia_solesmes_1961} %%the asterisk may need adjusting in the end.
\antiphona{english}{It is written: My house is the house of prayer for all nations but ye have made it a den of thieves. And he taught daily in the Temple}

\collectlateng

\texteparacol{\lettrine{P}{a}teant aures misericórdiæ tuæ Dómine précibus supplicántium:~† et ut peténtibus desideráta concédas,~* fac eos, quæ tibi sunt plácita postuláre.
Per Dóminum.}{english}
{\noindent Let thy merciful ears, O Lord, be open to the prayers of thy humble servants; and, that they may obtain their petitions, make them to ask such things as shall please thee.
Through our Lord.}

\bigtitle{X Sunday after Pentecost.}
\addcontentsline{toc}{section}{X Sunday after Pentecost} \phantomsection

\gscore[Ad Magnif.]{Ant. 8. G}{an_descendit_hic_solesmes_1961}
\antiphona{english}{This man went down to his house justified rather than the other, for every one that exalteth himself shall be abased, and he that humbleth himself shall be exalted}

\collectlateng

\texteparacol{\lettrine{D}{e}us, qui omnipoténtiam tuam parcéndo máxime et miserándo maniféstas:~† multíplica super nos misericórdiam tuam; ut ad tua promíssa curréntes,~* cæléstium bonórum fácias esse consórtes. Per Dóminum.}{english}
{\noindent O God, who declarest thine almighty power most chiefly in showing mercy and pity, mercifully grant unto us such a measure of thy grace, that we, running the way of thy commandments, may obtain thy gracious promises, and be made partakers of thy heavenly treasure.
Through our Lord.}

\bigtitle{XI Sunday after Pentecost.}
\addcontentsline{toc}{section}{XI Sunday after Pentecost} \phantomsection

\gscore[Ad Magnif.]{Ant. 5. a}{an_bene_omnia_solesmes_1961}
\antiphona{english}{He hath done all things well, he maketh both the deaf to hear and the dumb to speak}

\collectlateng

\texteparacol{\lettrine{O}{m}nípotens sempitéme Deus, qui abundántia pietátis tuæ, et mérita súpplicum excédis et vota:~† effúnde super nos misericórdiam tuam; ut dimíttas quæ consciéntia métuit,~* et adícias quod orátio non præsúmit. Per Dóminum.}{english}
{\noindent Almighty and everlasting God, who art always more ready to hear than we to pray, and art wont to give more than either we desire or deserve, pour down upon us the abundance of thy mercy, forgiving us those things whereof our conscience is afraid, and giving us those good things which we are not worthy to ask.
Through our Lord.}


\bigtitle{XII Sunday after Pentecost.}
\addcontentsline{toc}{section}{XII Sunday after Pentecost} \phantomsection

\gscore[Ad Magnif.]{Ant. 8. G}{an_homo_quidam_descendebat_solesmes_1961}
\antiphona{english}{A certain man went down from Jerusalem to Jericho, and fell among thieves, which stripped him of his raiment, and wounded him, and departed, leaving him half dead}

\collectlateng

\texteparacol{\lettrine{O}{m}nípotens et miséricors Deus, de cujus múnere venit, ut tibi a fidélibus tuis digne et laudabíliter serviátur:~† tríbue quǽsumus nobis,~* ut ad promissiónes tuas sine offensióne currámus. Per Dóminum.}{english}
{\noindent Almighty and merciful God, of whose only gift it cometh that thy faithful people do unto thee true and laudable service, grant, we beseech thee, that we may so faithfully serve thee in this life, that we fail not finally to attain thy heavenly promises.
Through our Lord.}

\bigtitle{XIII Sunday after Pentecost.}
\addcontentsline{toc}{section}{XIII Sunday after Pentecost} \phantomsection

\gscore[Ad Magnif.]{Ant. 1. D2}{an_unus_autem_ex_illis_solesmes_1961}
\antiphona{english}{And one of them, when he saw that he was healed, turned back, and with a loud voice glorified God, alleluia}

\collectlateng

\texteparacol{\lettrine{O}{m}nípotens sempitéme Deus, da nobis fídei, spei et caritátis augméntum:~† et, ut mereámur ássequi quod promíttis,~* fac nos amáre quod prǽcipis. Per Dóminum.}{english}{\noindent 
Almighty and everlasting God, give unto us the increase of faith, hope, and charity, and that we may worthily obtain that which Thou dost promise, make us to love that which Thou dost command. Through our Lord.}

\bigtitle{XIV Sunday after Pentecost.}
\addcontentsline{toc}{section}{XIV Sunday after Pentecost} \phantomsection

\gscore[Ad Magnif.]{Ant. 1. g}{an_quaerite_primum_solesmes_1961}
\antiphona{english}{Seek ye first the kingdom of God and his righteousness, and all these things shall be added unto you, alleluia}

\collectlateng

\texteparacol{\lettrine{C}{u}stódi Dómine, quǽsumus, Ecclésiam tuam propitiatióne perpétua:~† et quia sine te lábitur humána mortálitas,~* tuis semper auxíliis et abstrahátur a nóxiis, et ad salutária dirigátur. Per Dóminum.}{english}{\noindent 
Keep, we beseech thee, O Lord, thy Church with thy perpetual mercy, and because the frailty of man without thee cannot but fall, keep us ever by thy help from all things hurtful, and lead us to all things profitable to our salvation.
Through our Lord.}
 
\bigtitle{XV Sunday after Pentecost.}
\addcontentsline{toc}{section}{XV Sunday after Pentecost} \phantomsection

\gscore[Ad Magnif.]{Ant. 4. A}{an_propheta_magnus_solesmes}
\antiphona{english}{A great prophet is risen up among us and God hath visited his people}

\collectlateng

\texteparacol{\lettrine{E}{c}clésiam tuam Dómine miseratio continuáta mundet et múniat:~† et quia sine te non potest salva consístere,~* tuo semper múnere gubemátur. Per Dóminum.}
{english}{\noindent 
O Lord, we beseech thee, let thy continual pity cleanse and defend thy Church, and because it cannot continue in safety without thy succour, preserve it evermore by thy help and goodness. Through our Lord.}

\bigtitle{XVI Sunday after Pentecost.} 
\addcontentsline{toc}{section}{XVI Sunday after Pentecost} \phantomsection

\gscore[Ad Magnif.]{Ant. 7. a}{an_cum_vocatus_fueris_solesmes_1961}
\antiphona{english}{When thou art bidden of any man to a wedding, go and sit down in the lowest room, that he that bade thee may say: Friend, go up higher, then shalt thou have worship in the presence of them that sit at meat with thee, alleluia}

\collectlateng

\texteparacol{\lettrine{T}{u}a nos, quǽsumus Dómine, grátia semper et prævéniat et sequátur:~* ac bonis opéribus júgiter præstet esse inténtos.
Per Dóminum.}{english}{\noindent Lord, we pray thee, that thy grace may always prevent and follow us, and make us continually to be given to all good works.
Through our Lord.}

\bigtitle{XVII Sunday after Pentecost.} 
\addcontentsline{toc}{section}{XVII Sunday after Pentecost} \phantomsection

\gscore[Ad Magnif.]{Ant. 4. E}{an_quid_vobis_videtur_solesmes_1961}
\begin{enpars} \noindent Ant. What think ye of Christ? Whose Son is He? They say all unto him: The Son of David. Jesus saith unto them: How then doth David in spirit call him Lord, saying: The Lord said unto my Lord, Sit Thou at My right hand?\end{enpars}

\collectlateng

\texteparacol{\lettrine{D}{a}, quǽsumus Dómine, pópulo tuo diabólica vitáre contágia:~* et te solum Deum pura mente sectári. Per Dóminum.}{english}{\noindent 
O Lord, we beseech thee, grant thy people grace to withstand the temptations of the devil, and with pure hearts to follow thee the only God.
Through our Lord.}

\bigtitle{XVIII Sunday after Pentecost.}
\addcontentsline{toc}{section}{XVIII Sunday after Pentecost} \phantomsection 

\gscore[Ad Magnif.]{Ant. 4. E}{an_tulit_ergo_solesmes_1961}
\antiphona{english}{And immediately, he that had been sick of the palsy arose, and took up his bed whereon he lay, glorifying God and all the people, when they saw it, gave praise to God}

\collectlateng

\texteparacol{\lettrine{D}{i}rigat corda nostra, quǽsumus Dómine, tuæ miseratiónis operátio:~* quia tibi sine te placére non póssumus. Per Dóminum.}{english}{\noindent Mercifully grant, O Lord, that thine effectual goodness may in all things direct our hearts, forasmuch as without thee we are not able to please thee.
Through our Lord.}

\bigtitle{XIX Sunday after Pentecost.}
\addcontentsline{toc}{section}{XIX Sunday after Pentecost} \phantomsection

\gscore[Ad Magnif.]{Ant. 3. a}{an_intravit_autem_rex_solesmes_1961}
%\antiphona{english}{}
\begin{enpars} \noindent Ant. 
And the King came in to see the guests, and he saw there a man which had not on a wedding garment. And he saith unto him: Friend, how camest thou in hither, not having a wedding garment?\end{enpars}

\collectlateng

\texteparacol{\lettrine{O}{m}nípotens et miséricors Deus, univérsa nobis adversántia propitiátus exclúde:~† ut mente et córpore páriter expedíti,~* quæ tua sunt, líberis méntibus exsequámur. Per Dóminum.}{english}{\noindent O Almighty and most merciful God, of thy bountiful goodness keep us, we beseech thee, from all things that may hurt us; that we, being ready both in body and soul, may cheerfully accomplish those things that Thou wouldest have done.
Through our Lord.}

\bigtitle{XX Sunday after Pentecost.} 
\addcontentsline{toc}{section}{XX Sunday after Pentecost} \phantomsection

\gscore[Ad Magnif.]{Ant. 3. a}{an_cognovit_autem_solesmes_1961}
\antiphona{english}{So the father knew that it was at the same hour in which Jesus said “Thy son liveth.” And he himself believed, and his whole house}

\collectlateng

\texteparacol{\lettrine{L}{a}rgíre, quǽsumus Dómine, fidélibus tuis indulgéntiam placátus et pacem:~† ut páriter ab ómnibus mundéntur offénsis,~* et secúra tibi mente desérviant. Per Dóminum.}{english}{\noindent Grant, we beseech thee, O Lord, to thy faithful people pardon and peace, that they may be cleansed from all their sins, and serve thee with a quiet mind. Through our Lord.}

\bigtitle{XXI Sunday after Pentecost.}
\addcontentsline{toc}{section}{XXI Sunday after Pentecost} \phantomsection  

\gscore[Ad Magnif.]{Ant. 6. c}{an_serve_nequam_solesmes_1961}
\antiphona{english}{O thou wicked servant, I forgave thee all that debt, because thou desiredst me. Shouldest not thou also have had compassion on thy fellow-servant, even as I had pity on thee, alleluia}

\collectlateng

\texteparacol{\lettrine{F}{a}míliam tuam, quǽsumus Dómine, contínua pietáte custódi:~† ut a cun\-ctis adversitátibus te protegénte, sit líbera,~* et in bonis áctibus tuo nómini sit devóta. Per Dóminum.}{english}{\noindent O God, who seest that in our own weakness we do continually fall, make, in thy mercy, the examples of thy holy children a mean whereby to renew in us the love of thyself. Through our Lord.}

\bigtitle{XXII Sunday after Pentecost.}
\addcontentsline{toc}{section}{XXII Sunday after Pentecost} \phantomsection  

\gscore[Ad Magnif.]{Ant. 1. g}{an_reddite_ergo_solesmes_1961}
\antiphona{english}{Render therefore unto Caesar the things which are Caesar's, and unto God the things that are God's, alleluia}

\newpage
\collectlateng

\texteparacol{\lettrine{D}{e}us, refúgium nostrum et virtus:~† adésto piis Ecclésiæ tuæ précibus, au\-ctor ipse pietátis, et præsta;~* ut quod fidéliter pétimus, efficáciter consequámur. Per Dóminum.}{english}{\noindent O God, our refuge and strength, who art the author of all godliness, be ready, we beseech thee, to hear the devout prayers of thy Church, and grant that those things which we ask faithfully, we may obtain effectually.
Through our Lord.}

\rubrique{If the following Sunday is the last after Pentecost, in its place is placed XXIV, as below, and the office of this Sunday is placed on the preceding Saturday not impeded by a feast of IX lessons; otherwise on Friday or a similar day not impeded.}

\bigtitle{XXIII Sunday after Pentecost.}
\addcontentsline{toc}{section}{XXIII Sunday after Pentecost} \phantomsection  

\gscore[Ad Magnif.]{Ant. 1. f}{an_at_jesus_conversus_solesmes_1961}
\antiphona{english}{But Jesus turned him about, and when He saw her, He said: Daughter, be of good comfort; thy faith hath made thee whole, alleluia}

\collectlateng

\texteparacol{\lettrine{A}{b}sólve, quǽsumus Dómine, tuórum delícta populórum:~† ut a peccatórum néxibus, quæ pro nostra fragilitáte contráximus,~* tua benignitáte liberémur. Per Dóminum.}{english}{\noindent O Lord, we beseech thee, absolve thy people from their offences, that through thy bountiful goodness we may all be delivered from the bands of those sins, which by our frailty we have committed.
Through our Lord.}

\rubrique{The Sundays after Pentecost  cannot be fewer than XXIII, nor greater than XXVIII. But when there be more than XXIV, then after the XXIII Sunday, are resumed those Sundays which that year were left over after the Epiphany, with respect to the collect and the antiphons at the Magnificat, in this order:}

\rubrique{If the Sundays after Pentecost be XXV, the XXIV Sunday after Pentecost will be that which is VI after the Epiphany.}

\rubrique{If there be XXVI, the XXIV Sunday will be that which is V after the Epiphany, and XXV that which is VI.}

\rubrique{If there be XXVII, the XXIV Sunday will be that which is IV after the Epiphany, XXV will be V, and XXVI will be VI.}

\rubrique{If there be XXVIII Sundays, the XXIV Sunday will be that which is the III after the Epiphany, XXV will be IV, XXVI will be V, and XXVII will be VI.}

\rubrique{The last place is always taken by that Sunday which follows in order as XXIV.}

\bigtitle{III Sunday which was left over after the Epiphany.}

\rubrique{At the Magníficat, ant. \normaltext{Dómine, si tu vis, \pageref{M-3_sunday_pe}.}}

\rubrique{Collect \normaltext{Omnípotens sempitérne Deus.}}

\bigtitle{IV Sunday which was left over after the Epiphany.}

\rubrique{At the Magníficat, ant. \normaltext{Dómine, salva nos, \pageref{M-4_sunday_pe}.}}

\rubrique{Collect \normaltext{Deus, qui nos in tantis perículis.}}

\bigtitle{V Sunday which was left over after the Epiphany.}

\rubrique{At the Magníficat, ant. \normaltext{Collígite primum, \pageref{M-5_sunday_pe}.}}

\rubrique{Collect \normaltext{Famíliam tuam, quǽsumus Dómine.}}

\bigtitle{VI Sunday which was left over after the Epiphany.}

\rubrique{At the Magníficat, ant. \normaltext{Símile est regnum cælórum ferménto, \pageref{M-6_sunday_pe}.}}

\rubrique{Collect \normaltext{Præsta, quǽsumus omnípotens Deus.}}

\bigtitle{XXIV Sunday and the last after Pentecost.}
\addcontentsline{toc}{section}{XXV Sunday after Pentecost} \phantomsection    

\gscore[Ad Magnif.]{Ant. 1. f}{an_amen_dico_vobis_quia_non_solesmes_1961}
\antiphona{english}{Amen, I say unto you, this generation shall not pass till all these things be fulfilled; heaven and earth shall pass away, but my words shall not pass away, saith the Lord}

\collectlateng

\texteparacol{\lettrine{E}{x}cita, quǽsumus Dómine, tuórum fidélium voluntátes:~† ut divíni óperis fru\-ctum propénsius exsequéntes,~* pietátis tuæ remédia majóra percípiant.
Per Dóminum.}{english}{\noindent 
Stir up, we beseech thee, Lord, the wills of thy faithful people, that they, plenteously bringing forth the fruit of good works, may of thee be plenteously rewarded.
Through our Lord.}

\end{document}